%! Author = lili
%! Date = 6/15/26

\section{The Forward-Backward Technique}
\skript{The following problem frequently arises as a sub-problem of more advanced alignment methods, that is
why we discuss it beforehand in detail:}

\begin{prob}[Advanced Alignment Problem]
\skript{Given two sequences $x$ and $y$ of lengths $m$ and $n$, respectively, find their optimal global
alignment under the condition that the alignment path passes through a given cell $(i_0, j_0)$ of the
alignment matrix.}
\end{prob}

\skript{This is not a new problem at all! In fact, we can decompose it into two standard global alignment
problems: Aligning the prefixes $x[1\ldots i_0]$ and $y[1\ldots j_0]$, and aligning the suffixes
    $x[i_0+1\ldots m]$ and $y[j_0+1\ldots n]$. By combining the results of these two sub-problems in an
    appropriate way, we obtain the minimally attainable total cost of all alignment paths that go through
    $(i_0, j_0)$.

    \paragraph{Total cost matrix.}
    In fact, it is even possible to solve this problem for every point $(i,j)$ simultaneously. The
    solution is called the \emph{forward-backward technique}. In addition to the usual alignment matrix
    with alignment costs $D(i,j)$ (which record the minimally attainable cost for a prefix alignment), we
    additionally define the \emph{reverse alignment matrix}: We exchange initial and final cell and
    reverse the direction of all computations. Effectively, we are thus globally aligning the reversed
    sequences. Hence, the associated alignment matrix $D^{\mathrm{rev}} = D^{\mathrm{rev}}(i,j)$ records
    the best possible costs for the suffix alignments.}

\definSB{total_cost_matrix}

\begin{notte}
    \skript{Note that $T(0,0) = T(m,n)$ is the optimal global alignment cost $d(x,y)$ (since all paths, in
    particular the optimal ones, pass through $(0,0)$ and $(m,n)$), and so $T(i,j)\ge d(x,y)$ for all
        $(i,j)$.}
\end{notte}

\skript{
    The computation of the total cost matrix is particularly simple for additive alignment costs. Here
    \[
        T(i,j) = D(i,j) + D^{\mathrm{rev}}(i,j).
    \]
    Clearly matrices $D$, $D^{\mathrm{rev}}$, and thus $T$ can be computed in $O(mn)$ time.

    For affine gap costs the computation of the total cost matrix is slightly more difficult: Here, the
    cost of a concatenated alignment is not always the sum of the costs of the individual alignments
    because gaps might be merged, so that the gap initiation penalty that is imposed twice in the two
    separate alignments must be counted only once in the concatenated alignment. However, since the
    efficient computation of affine gap costs with Gotoh's algorithm uses the two history matrices $V$ and
    $H$, and we know that by definition the costs stored in these matrices refer to those alignments that
    end with gaps, the total cost matrix for affine gap costs can also be computed in quadratic time if
    the history matrices $V$ and $H$ and their reverse counterparts $V^{\mathrm{rev}}$ and
    $H^{\mathrm{rev}}$ are given:
    \[
        T(i,j) = \min
        \begin{cases}
            D(i,j) + D^{\mathrm{rev}}(i,j)\\
            V(i,j) + V^{\mathrm{rev}}(i,j) - d + e\\
            H(i,j) + H^{\mathrm{rev}}(i,j) - d + e
        \end{cases}
    \]
    where $d$ is the gap open cost and $e$ is the gap extension cost.

    \paragraph{Additional cost matrix.}
    For every $(i,j)$, we can now obtain the \emph{additional cost} $C(i,j)$ for passing through $(i,j)$,
    compared to staying on an optimal alignment path.}

\definSB{additional_cost_matrix}

% todo By definition, if $(i,j)$ is on the path of an optimal alignment, this value is zero, otherwise it is
%positive. %; see Figure~6.1.

\skript{By definition, if $(i,j)$ is on the path of an optimal alignment, this value is zero, otherwise it is
positive. %; see Figure~6.1.

Again, for any additive alignment cost, where
    $\mathrm{cost}(A++B) = \mathrm{cost}(A) + \mathrm{cost}(B)$, we have that
    $C(i,j) = T(i,j) - d(x,y) = D(i,j) + D^{\mathrm{rev}}(i,j) - d(x,y)$. Since $T$ and $d(x,y)$ can be
    computed in $O(mn)$ time, the additional cost matrix $C$ can be computed in $O(mn)$ time, too.
}

\begin{bsp}{Aus Skript}

    Consider the sequences $x = \mathtt{CT}$ and $y = \mathtt{AGT}$. For unit cost, the edit matrix $D$ and
    reverse edit matrix $D^{\mathrm{rev}}$, which is drawn in reverse direction, are:
    \[
        D:\quad
        \begin{array}{c|cccc}
            & \epsilon & \mathtt{A} & \mathtt{G} & \mathtt{T}\\\hline
            \epsilon & 0 & 1 & 2 & 3\\
            \mathtt{C} & 1 & 1 & 2 & 3\\
            \mathtt{T} & 2 & 2 & 2 & 2
        \end{array}
        \qquad\qquad
        D^{\mathrm{rev}}:\quad
        \begin{array}{c|cccc}
            & \mathtt{A} & \mathtt{G} & \mathtt{T} & \epsilon\\\hline
            \mathtt{C} & 2 & 1 & 1 & 2\\
            \mathtt{T} & 2 & 1 & 0 & 1\\
            \epsilon & 3 & 2 & 1 & 0
        \end{array}
    \]
    This results in the following additional cost matrix:
    \[
        C:\quad
        \begin{array}{c|cccc}
            & \epsilon & \mathtt{A} & \mathtt{G} & \mathtt{T}\\\hline
            \epsilon & 0 & 0 & 1 & 3\\
            \mathtt{C} & 1 & 0 & 0 & 2\\
            \mathtt{T} & 3 & 2 & 1 & 0
        \end{array}
    \]
    We see that the paths of the two optimal alignments
    $A^{\mathrm{opt}}_1 = \left(\begin{smallmatrix}\mathtt{C}&\mathtt{-}&\mathtt{T}\\
    \mathtt{A}&\mathtt{G}&\mathtt{T}\end{smallmatrix}\right)$ and
    $A^{\mathrm{opt}}_2 = \left(\begin{smallmatrix}\mathtt{-}&\mathtt{C}&\mathtt{T}\\
    \mathtt{A}&\mathtt{G}&\mathtt{T}\end{smallmatrix}\right)$ correspond to the two paths of $0$-entries in
    $C$.
\end{bsp}

\skript{\paragraph{Applications.}
The forward-backward technique is used, for example, in the following problems:
    \begin{itemize}
        \item linear-space alignment (Hirschberg technique, see Section~6.2),
        \item discovering regions of slightly sub-optimal alignments (those cells where the score gap, resp.\
        additional cost, remains below a small tolerance parameter),
        \item computing regions for the Carrillo-Lipman heuristic for multiple sequence alignment (see
        Section~8.2),
        \item finding cut-points for divide-and-conquer multiple alignment (see Section~10.1).
    \end{itemize}}

\section{Hirschberg’s Algorithm}
TODO